\documentclass[12pt,a4paper]{article}
\usepackage[utf8]{inputenc}
\usepackage[T1]{fontenc}
\usepackage{amsmath,amssymb,amsthm,mathtools}
\usepackage{geometry}
\geometry{left=1.0cm,right=1.0cm,top=1.5cm,bottom=2.0cm}
\usepackage{booktabs}
\usepackage{tabularx}
\usepackage{array}
\usepackage{hyperref}
\usepackage{url}
\usepackage{xcolor}
\usepackage{titlesec}
\usepackage{fancyhdr}
\usepackage{microtype}
\usepackage{enumitem}
\usepackage{slashed}
\usepackage{physics}
\usepackage{float}

% YUCT standard fonts
\usepackage{fontspec}
\setmainfont{Calibri}
\setsansfont{Segoe UI}[
BoldFont = Segoe UI Bold,
Ligatures = TeX
]

% YUCT color scheme
\definecolor{skyblue}{RGB}{0,102,204}
\definecolor{darkblue}{RGB}{0,0,139}
\definecolor{gold}{RGB}{255,215,0}

% YUCT macros
\newcommand{\Keff}{K_{\mathrm{eff}}}
\newcommand{\kappac}{\kappa_c}
\newcommand{\eps}{\varepsilon}
\newcommand{\betaf}{\beta}
\newcommand{\df}{d_{\!f}}
\newcommand{\Sodd}{S_{\mathrm{odd}}}
\newcommand{\Seven}{S_{\mathrm{even}}}
\newcommand{\Mpl}{M_{\mathrm{Pl}}}
\newcommand{\Lpl}{l_{\mathrm{Pl}}}
\newcommand{\tpl}{t_{\mathrm{Pl}}}
\newcommand{\Lmin}{\ell_{\min}}
\newcommand{\LUV}{\Lambda_{\mathrm{UV}}}

% Theorem environments
\newtheorem{theorem}{Theorem}[section]
\newtheorem{proposition}[theorem]{Proposition}
\newtheorem{definition}[theorem]{Definition}
\theoremstyle{remark}
\newtheorem{remark}[theorem]{Remark}

% Section formatting
\titleformat{\section}{\LARGE\bfseries\color{darkblue}}{\thesection}{1em}{}
\titleformat{\subsection}{\Large\bfseries\color{darkblue}}{\thesubsection}{1em}{}

% Headers and footers
\fancypagestyle{firstpage}{%
	\fancyhf{}%
	\renewcommand{\headrulewidth}{0pt}%
	\renewcommand{\footrulewidth}{0pt}%
	\fancyfoot[C]{%
		\begin{minipage}[t]{\textwidth}
			\centering
			\textcolor{gold}{\rule{\textwidth}{1.6pt}}\\[5pt]
			{\small \textsuperscript{1}Yakushev Research, \url{https://yuct.org/}} \hfill
			{\small YPSDC Protocol: \url{https://ypsdc.com/}}\\[-2pt]
			\textcolor{gold}{\rule{\textwidth}{1.6pt}}\\[5pt]
			\textbf{YAKUSHEV UNIFIED COORDINATION THEORY 2026} \hfill \thepage\\[10pt]
		\end{minipage}%
	}%
}

\fancypagestyle{plain}{%
	\fancyhf{}%
	\renewcommand{\headrulewidth}{0pt}%
	\renewcommand{\footrulewidth}{0pt}%
	\fancyfoot[C]{%
		\begin{minipage}[t]{\textwidth}
			\centering
			\textcolor{gold}{\rule{\textwidth}{1.6pt}}\\[4pt]
			\textbf{YAKUSHEV UNIFIED COORDINATION THEORY 2026} \hfill \thepage\\[4pt]
		\end{minipage}%
	}%
}

\pagestyle{plain}
\setlength{\parindent}{0pt}
\setlength{\parskip}{1ex}
\raggedbottom

\hypersetup{
	colorlinks=true,
	linkcolor=blue,
	citecolor=blue,
	urlcolor=blue,
}

% Custom header
\newcommand{\makeheader}{%
	\noindent
	\begin{minipage}{0.17\textwidth}
		\raggedright
		{\fontspec{Segoe UI}[BoldFont=Segoe UI Bold]\bfseries\color{skyblue}\fontsize{46}{46}\selectfont YUCT}%
	\end{minipage}%
	\hspace{30pt}
	\begin{minipage}{0.32\textwidth}
		\raggedright
		{\sffamily\fontsize{11}{13}\selectfont YAKUSHEV\\ UNIFIED\\ COORDINATION THEORY}%
	\end{minipage}%
	\hfill
	\begin{minipage}{0.38\textwidth}
		\raggedleft
		{\sffamily\bfseries Appendix UVD}\\ 
		{\sffamily Version 1.0 / June 2026}\\ 
		{\sffamily\footnotesize \url{https://doi.org/10.5281/zenodo.18444598}}%
	\end{minipage}
	\par\vspace{5pt}
	\noindent\textcolor{gold}{\rule{\textwidth}{1.6pt}}
	\par\vspace{0pt}
}

\begin{document}
	
	% Title page
	\thispagestyle{firstpage}
	\makeheader
	
	\vspace{0.5cm}
	{\LARGE\bfseries Appendix UVD: Resolution of Ultraviolet Divergences in YUCT \\
		\Large Finite Quantum Field Theory from the Coordination Form Factor and the Physical Scale of the Coordination Network \par}
	\vspace{0.3cm}
	
	{\large Alexey V. Yakushev\textsuperscript{1}} \\
	\vspace{0.2cm}
	
	{\color{darkblue}\textbf{\LARGE Abstract}} \par
	{\large
		\noindent
		The Yakushev Unified Coordination Theory (YUCT) proposes that ultraviolet divergences in quantum field theory are eliminated by a coordination form factor \\ \(F(q^2) = \exp[-(q^2/\LUV^2)^{2/3}]\) arising from the finite information capacity of the vacuum.  In the original formulation, the UV scale \(\LUV\) was identified with the Planck mass, \(\LUV = 3\pi\Mpl \approx 1.15\times 10^{20}\)~GeV, leading to an enormous finite radiative correction to the Higgs mass.  We correct a calculational error in the constant \(C_0\) appearing in the scalar mass correction: the correct value is \(C_0 = 3\sqrt{\pi}/(8\sqrt{2}) \approx 0.46999\), replacing the previously reported \(C_0 \approx 0.626\).  With this correction, the radiative contribution to the Higgs mass from a Planck-scale network is \(\delta m \approx 1.57\times 10^{18}\)~GeV, which reintroduces a severe fine-tuning problem unless the bare mass is precisely cancelled.  We examine two logical alternatives: (i) the YUCT mass formula \(m = \Mpl(2/3)^L\xi_f\) already determines the physical mass, and the radiative correction is absorbed by a finite counterterm -- a technically consistent but aesthetically unsatisfying solution; (ii) the fundamental Lagrangian at the Planck scale contains no bare mass term (scale invariance), and the entire Higgs mass is generated by the radiative correction.  In this second scenario, self-consistency demands that the UV scale be \(\LUV \approx 8.96\)~TeV, a value accessible to the LHC.  The coordination form factor then predicts observable deviations in high-energy scattering cross sections at the TeV scale.  We present both scenarios, their mathematical justification, and their experimental consequences, with the aim of subjecting YUCT to rigorous empirical testing.
	}
	
	\vspace{0.2cm}
	\noindent
	{\color{darkblue}\textbf{Keywords:}} YUCT, ultraviolet divergences, coordination form factor, hierarchy problem, Higgs mass, scale invariance, nonlocal quantum field theory, LHC phenomenology.
	\vspace{0.1cm}
	\noindent
	{\color{darkblue}\textbf{Citation:}} Yakushev AV. Appendix UVD: Resolution of Ultraviolet Divergences in YUCT. 2026. DOI: \url{https://doi.org/10.5281/zenodo.18444598} (this volume).
	
	\vspace{0.5cm}
	
	% ============================================================
	\section{Introduction}
	\label{sec:intro}
	
	The Yakushev Unified Coordination Theory (YUCT) offers an elegant resolution to the ultraviolet divergences of quantum field theory (QFT).  Instead of introducing an arbitrary momentum cut-off or invoking infinite renormalisation, YUCT posits that the physical vacuum possesses a finite ``bandwidth'' for transmitting quantum information.  This bandwidth is quantified by a \textbf{coordination form factor} \(F(q^2)\) that exponentially suppresses virtual momenta above a characteristic scale \(\LUV\).  The mathematical structure of this form factor is dictated by the universal error law of YUCT, \(\varepsilon \propto K_{\mathrm{eff}}^{-\beta}\) with \(\beta = 2/3\).
	
	In previous versions of this appendix, we identified \(\LUV\) with the Planck mass, \(\LUV = 3\pi\Mpl \approx 1.15\times 10^{20}\)~GeV, a scale derived from the minimal length \(\Lmin = \frac{2}{3}\Lpl\) dictated by the primary algebraic loop of YUCT.  With this identification, all loop integrals become manifestly finite: logarithmic divergences are replaced by \(\ln(\LUV/m)\), and power-law divergences vanish identically.
	
	However, a careful recalculation of the scalar mass correction reveals a computational error in the constant \(C_0\) that governs the magnitude of the finite remnant.  The correct value of \(C_0\) is smaller than previously reported, reducing the radiative correction but still leaving it many orders of magnitude larger than the observed Higgs mass.  This forces a critical re-examination of the physical scale \(\LUV\).
	
	This appendix has two objectives.  First, we present the corrected mathematical derivations in full detail, ensuring that all expressions are accurate and reproducible.  Second, we explore the physical implications of the corrected formulae, including a radical proposal: that the fundamental scale of the coordination network may not be the Planck mass at all, but a much lower scale, potentially within reach of current collider experiments.
	
	% ============================================================
	\section{The Coordination Form Factor and the UV Scale}
	\label{sec:formfactor}
	
	The YUCT coordination form factor is
	\begin{equation}
		F(q^2) = \exp\!\Bigl[-\Bigl(\frac{q^2}{\LUV^2}\Bigr)^{\!\beta}\,\Bigr], \qquad \beta = \frac{2}{3}.
		\label{eq:F}
	\end{equation}
	All Standard Model propagators are multiplied by \(F(p^2/\LUV^2)\).  The scale \(\LUV\) is originally derived from the minimal length \(\Lmin = \frac{2}{3}\Lpl\) via the Compton relation \(\LUV = 2\pi\hbar c/\Lmin = 3\pi\Mpl\).
	
	This form factor renders all loop integrals absolutely convergent in the ultraviolet.  The proof is straightforward: for any polynomial \(P(k)\) in the loop momenta, the product \(P(k) \cdot \prod_i F(k_i^2/\LUV^2)\) decays faster than any inverse power of \(|k|\) as \(|k| \to \infty\), because the exponential dominates.
	
	% ============================================================
	\section{Electron Self-Energy: A Logarithmic Divergence Resolved}
	\label{sec:electron}
	
	The one-loop electron self-energy in QED is logarithmically divergent in conventional QFT.  With the YUCT form factor, the Euclidean integral becomes
	\begin{equation}
		\Sigma_{\mathrm{YUCT}} \propto \int_0^{\infty} \frac{dk_E}{k_E}\; \exp\!\Bigl[-2\Bigl(\frac{k_E^2}{\LUV^2}\Bigr)^{\!2/3}\,\Bigr].
		\label{eq:elec_int}
	\end{equation}
	The substitution \(t = (k_E^2/\LUV^2)^{2/3}\) yields \(dk_E/k_E = \frac{3}{4} dt/t\), and the integral evaluates to
	\begin{equation}
		I_{\mathrm{log}} = \frac{3}{4} \int_{t_{\min}}^{\infty} \frac{dt}{t}\, e^{-2t} \approx \ln\frac{\LUV}{m_e},
		\label{eq:elec_result}
	\end{equation}
	where \(t_{\min} \sim (m_e^2/\LUV^2)^{2/3}\).  The logarithmic divergence is replaced by the finite logarithm \(\ln(\LUV/m_e) \approx 53.8\).  The calculation is elementary and requires no correction.
	
	% ============================================================
	\section{Scalar Mass Correction: Correcting the Constant \(C_0\)}
	\label{sec:scalar}
	
	\subsection{The Integral and Its Evaluation}
	
	The one-loop scalar mass correction in \(\phi^4\) theory, after Wick rotation and insertion of the form factor, is
	\begin{equation}
		\delta m^2_{\mathrm{YUCT}} = \frac{\lambda}{16\pi^2} \int_0^{\infty} \frac{k_E^3\, dk_E}{k_E^2 + m^2}\; F^2\!\Bigl(\frac{k_E^2}{\LUV^2}\Bigr), \qquad F^2(x) = e^{-2x^{2/3}}.
		\label{eq:scalar_YUCT}
	\end{equation}
	
	Introducing the dimensionless variable \(t = k_E^2/\LUV^2\) and expanding for \(m \ll \LUV\), the leading term is
	\begin{equation}
		\delta m^2_{\mathrm{YUCT}} = \frac{\lambda}{32\pi^2}\,\LUV^2 \int_0^{\infty} e^{-2t^{2/3}} dt + \mathcal{O}(m^2/\LUV^2).
		\label{eq:scalar_lead}
	\end{equation}
	
	The integral \(C_0 \equiv \int_0^{\infty} e^{-2t^{2/3}} dt\) was previously reported as \(C_0 \approx 0.626\) (version 3.1).  This value is incorrect.  We now evaluate it exactly.
	
	\subsection{Exact Evaluation of \(C_0\)}
	
	\begin{theorem}[Exact Value of \(C_0\)]
		\label{thm:C0}
		\[
		C_0 = \int_0^{\infty} e^{-2t^{2/3}} dt = \frac{3\sqrt{\pi}}{8\sqrt{2}} \approx 0.46999.
		\]
	\end{theorem}
	
	\begin{proof}
		Substitute \(u = 2t^{2/3}\).  Then \(t = (u/2)^{3/2}\) and
		\[
		dt = \frac{3}{2} (u/2)^{1/2} \cdot \frac{1}{2} du = \frac{3}{4\sqrt{2}} u^{1/2} du.
		\]
		Hence
		\[
		C_0 = \int_0^{\infty} e^{-u} \cdot \frac{3}{4\sqrt{2}} u^{1/2} du = \frac{3}{4\sqrt{2}} \int_0^{\infty} u^{1/2} e^{-u} du = \frac{3}{4\sqrt{2}} \, \Gamma(3/2).
		\]
		Using \(\Gamma(3/2) = \frac{\sqrt{\pi}}{2}\), we obtain
		\[
		C_0 = \frac{3}{4\sqrt{2}} \cdot \frac{\sqrt{\pi}}{2} = \frac{3\sqrt{\pi}}{8\sqrt{2}} \approx 0.46999.
		\]
	\end{proof}
	
	\subsection{Corrected Scalar Mass Correction}
	
	With the correct value of \(C_0\), the leading scalar mass correction is
	\begin{equation}
		\boxed{\delta m^2_{\mathrm{YUCT}} = \frac{\lambda}{32\pi^2}\,\LUV^2 \cdot \frac{3\sqrt{\pi}}{8\sqrt{2}} + \mathcal{O}(m^2/\LUV^2)}.
		\label{eq:scalar_corrected}
	\end{equation}
	
	For the Higgs boson, \(\lambda = m_H^2/(2v^2) \approx 0.13\) (where \(v \approx 246\)~GeV).  Using the Planck-scale identification \(\LUV = 3\pi\Mpl \approx 1.15\times 10^{20}\)~GeV:
	\begin{align}
		\delta m^2_H &\approx \frac{0.13}{32\pi^2} \cdot (1.15\times 10^{20})^2 \cdot 0.46999 \notag\\
		&\approx 2.47 \times 10^{36}\ \mathrm{GeV}^2, \notag\\
		\delta m_H &\approx 1.57 \times 10^{18}\ \mathrm{GeV}.
		\label{eq:planck_correction}
	\end{align}
	
	This value exceeds the observed Higgs mass (\(m_H^{\mathrm{exp}} \approx 125\)~GeV) by a factor of \(1.26 \times 10^{16}\).  If the bare mass in the Lagrangian is zero, the physical mass is determined entirely by this radiative correction, and the predicted value is clearly incompatible with observation.
	
	% ============================================================
	\section{Physical Implications: Two Scenarios}
	\label{sec:scenarios}
	
	The corrected calculation forces us to confront the hierarchy problem directly.  We examine two logically consistent alternatives.
	
	\subsection{Scenario A: Planck-Scale Network with Finite Counterterms}
	
	In this scenario, YUCT retains the Planck-scale identification \(\LUV = 3\pi\Mpl\).  The physical mass of any particle is given by the YUCT mass formula
	\begin{equation}
		m_{\mathrm{phys}} = \Mpl \left(\frac{2}{3}\right)^L \xi_f,
		\label{eq:mass_formula}
	\end{equation}
	where \(L\) is the coordination depth and \(\xi_f\) is a resonance factor (see Appendix~J for a detailed discussion, including the origin of \(\xi_f\) from the topology of the compact fibre \(F_{15}\)).  The bare mass in the Lagrangian is a free parameter, chosen to cancel the radiative correction and reproduce the observed mass.
	
	Mathematically, this is perfectly consistent: the theory is finite, the required counterterm is also finite, and no infinity is ever encountered.  The YUCT mass formula, Eq.~(\ref{eq:mass_formula}), provides the renormalisation condition that fixes the counterterm.
	
	Physically, however, this scenario requires the bare mass and the radiative correction to cancel to one part in \(10^{16}\) -- a fine-tuning of precisely the kind that the hierarchy problem was meant to avoid.  YUCT does not explain why the bare mass should be chosen to achieve this cancellation; it merely replaces the usual infinite subtraction with an enormous finite one.
	
	\subsection{Scenario B: Scale-Invariant Origin and a TeV-Scale Network}
	
	An alternative, aesthetically more appealing resolution is to abandon the Planck-scale identification of \(\LUV\).  If the fundamental theory at the Planck scale possesses no bare mass terms (i.e., it is scale-invariant), then all masses are generated radiatively by the coordination form factor.  In this case, the physical Higgs mass \emph{is} the radiative correction:
	\begin{equation}
		m_H^2 = \delta m^2_{\mathrm{YUCT}} = \frac{\lambda}{32\pi^2}\,\LUV^2 \cdot C_0.
		\label{eq:scale_inv}
	\end{equation}
	
	Solving for \(\LUV\):
	\begin{equation}
		\LUV = \sqrt{\frac{32\pi^2\, m_H^2}{\lambda\, C_0}}.
		\label{eq:LUV_solve}
	\end{equation}
	
	Substituting the corrected \(C_0 = 3\sqrt{\pi}/(8\sqrt{2}) \approx 0.46999\), \(\lambda = 0.13\), and \(m_H = 125\)~GeV, we obtain
	\begin{equation}
		\boxed{\LUV \approx 8.96\ \mathrm{TeV}}.
		\label{eq:LUV_TeV}
	\end{equation}
	
	This value is a prediction of the theory, not a free parameter.  It places the fundamental coordination scale squarely in the TeV range, making it experimentally accessible.
	
	\subsection{Physical Interpretation of Scenario B}
	
	If \(\LUV \approx 9\)~TeV, the coordination network ``crystallises'' at the electroweak scale, not at the Planck scale.  The form factor \(F(q^2) = \exp[-(q^2/\LUV^2)^{2/3}]\) begins to suppress virtual momenta above a few TeV.  This has immediate observable consequences:
	
	\begin{enumerate}[label=(\arabic*)]
		\item \textbf{Drell-Yan and dijet cross sections} at the LHC should show a deficit of events at high invariant masses compared to Standard Model predictions.  The suppression becomes appreciable for \(m_{jj} \gtrsim 5\)~TeV.
		\item \textbf{Higgs pair production} and vector boson scattering should exhibit anomalous energy dependence as the form factor modifies the Higgs self-coupling.
		\item \textbf{Deviations in electroweak precision observables} may be expected, though they are likely small since the LEP measurements probe scales well below \(\LUV\).
	\end{enumerate}
	
	Crucially, the minimal length in this scenario is \(\Lmin = 2\pi\hbar c/\LUV \approx 1.38\times 10^{-19}\)~m, which is much larger than the Planck length (\(\sim 10^{-35}\)~m).  This implies that space-time granularity is not a quantum gravitational effect but an electroweak-scale phenomenon, potentially related to the Higgs mechanism itself.
	
	\subsection{The Fractal Bridge Between the Electroweak Scale and the Cosmological Constant}
	\label{subsec:bridge}
	
	If Scenario B is correct and the fundamental coordination scale is 
	\(\LUV \approx 8.96\)~TeV, a profound connection to cosmology emerges.  
	In YUCT, the vacuum is not a single continuous medium but a fractal 
	hierarchy of coordination layers.  The scale \(\LUV\) characterises the 
	layer at which the Standard Model fields experience the coordination 
	form factor.  However, deeper (infrared) layers exist, corresponding to 
	progressively lower energy scales.  The cosmological constant, i.e., 
	the observed dark energy density, is generated by the deepest infrared 
	layer, whose scale \(\Lambda_{\mathrm{IR}}\) is set by the requirement 
	that the zero‑point energy of all quantum fields, when integrated over 
	the full fractal hierarchy, reproduces the observed value 
	\(\rho_{\Lambda} \approx 2.51 \times 10^{-47}\)~GeV\(^4\) (see 
	Appendix~M for the complete derivation).
	
	The relationship between the UV scale \(\LUV\) and the IR scale 
	\(\Lambda_{\mathrm{IR}}\) is governed by the fundamental scaling quantum 
	of YUCT:
	\begin{equation}
		q = \left(\frac{3}{2}\right)^{1/3} \approx 1.1447.
	\end{equation}
	The two scales are separated by an integer number \(N\) of 
	coordination layers:
	\begin{equation}
		\LUV = \Lambda_{\mathrm{IR}} \cdot q^{N}.
		\label{eq:bridge}
	\end{equation}
	
	From the observed cosmological constant, one extracts 
	\(\Lambda_{\mathrm{IR}} \approx 0.0032\)~eV (see Appendix~M, 
	Eq.~(M.47)).  With \(\LUV \approx 8.96 \times 10^{12}\)~eV, we obtain
	\begin{equation}
		N = \frac{\ln(\LUV / \Lambda_{\mathrm{IR}})}{\ln q}
		= \frac{\ln(8.96 \times 10^{12} / 0.0032)}{\frac{1}{3}\ln(3/2)}
		\approx \frac{35.564}{0.135155} \approx 263.1.
		\label{eq:N_layers}
	\end{equation}
	
	Rounding to the nearest half‑integer (as required by the algebraic 
	loop of YUCT for fermionic layers), we obtain the fundamental 
	structural number
	\begin{equation}
		\boxed{N = 263.0}.
	\end{equation}
	
	This is a remarkable result.  The electroweak scale (the mass of the 
	Higgs boson) and the cosmological constant are connected by exactly 
	\(263\) discrete steps of the coordination hierarchy, without any 
	free parameters.  The observed smallness of the cosmological constant 
	is not a mystery; it simply reflects the enormous depth of the 
	coordination network that separates particle physics from cosmology.
	
	\textbf{Interpretation.}  Each layer of the hierarchy acts as a 
	``scale transformer'': the zero‑point energy of a given layer is 
	suppressed by the coordination form factor of the next, deeper layer.  
	The total vacuum energy is the sum of contributions from all layers, 
	each smaller than the previous by a factor of \(\beta = 2/3\).  
	The geometric series converges to a finite value, which is precisely 
	the observed dark energy density (see Appendix~M for the detailed 
	summation).  The integer \(N \approx 263\) is the number of layers 
	that fit between the electroweak scale and the neutrino mass scale 
	(the deepest physically accessible layer).  That this number emerges 
	as a half‑integer is a non‑trivial consistency check of the YUCT 
	framework, lending further support to Scenario~B.
	
	\subsection{Comparison of the Two Scenarios}
	
	\begin{table}[H]
		\centering
		\caption{Comparison of Scenario A and Scenario B.}
		\label{tab:scenarios}
		\small
		\begin{tabular}{@{}lcc@{}}
			\toprule
			\textbf{Property} & \textbf{Scenario A} & \textbf{Scenario B} \\
			\midrule
			UV scale \(\LUV\) & \(3\pi\Mpl \approx 10^{20}\)~GeV & \(\approx 8.96\)~TeV \\
			Minimal length \(\Lmin\) & \(\frac{2}{3}\Lpl \approx 10^{-35}\)~m & \(\approx 10^{-19}\)~m \\
			Bare mass in Lagrangian & Non-zero (fine-tuned) & Zero (scale invariance) \\
			Radiative correction \(\delta m_H\) & \(10^{18}\)~GeV & 125~GeV \\
			Experimental accessibility & Planck scale (inaccessible) & LHC energies (accessible) \\
			Fine-tuning required & Yes (1 part in \(10^{16}\)) & No \\
			\bottomrule
		\end{tabular}
	\end{table}
	
	% ============================================================
	\section{Falsifiable Predictions of Scenario B}
	\label{sec:predictions}
	
	If Scenario B is correct, the following quantitative predictions can be tested at the LHC and future colliders:
	
	\begin{enumerate}[label=(\arabic*)]
		\item \textbf{Drell-Yan cross section suppression.}  The ratio of measured to Standard Model Drell-Yan cross sections at dilepton invariant mass \(m_{\ell\ell}\) should follow
		\[
		R(m_{\ell\ell}) \approx \exp\!\Bigl[-2\Bigl(\frac{m_{\ell\ell}^2}{\LUV^2}\Bigr)^{\!2/3}\,\Bigr],
		\]
		with \(\LUV \approx 8.96\)~TeV.  A 10\% suppression is predicted at \(m_{\ell\ell} \approx 5.4\)~TeV, growing to a factor of \(\sim 2\) at \(m_{\ell\ell} \approx 10\)~TeV.
		
		\item \textbf{Higgs pair production.}  The form factor modifies the Higgs self-coupling at high momentum transfer.  The di-Higgs invariant mass distribution should deviate from the Standard Model for \(m_{HH} \gtrsim 3\)~TeV.
		
		\item \textbf{Fermion masses.}  If all fermion masses are also generated radiatively, their Yukawa couplings must be chosen such that the radiative corrections reproduce the observed mass spectrum.  This imposes a self-consistency condition relating the fermion coordination depths \(L_f\) and the resonance factors \(\xi_f\) to the radiative corrections.  A detailed analysis is deferred to a future work.
	\end{enumerate}
	
	\subsection{A Note on the Fine-Structure Constant}
	
	The derivation of \(\alpha^{-1} \approx 137.036\) from the topological invariant \(N_{\mathrm{gauge}} = 12\) (Appendix~G) is independent of the value of \(\LUV\).  The gauge couplings run logarithmically; the initial values at \(\LUV\) are fixed by the topology of \(F_{18}\).  In Scenario A, \(\LUV \sim \Mpl\), and the logarithmic running from the Planck scale to the electroweak scale yields \(\alpha^{-1} \approx 137.036\).  In Scenario B, \(\LUV \sim 9\)~TeV, and the running is over a much shorter interval.  The topological prediction for \(\alpha^{-1}\) would then apply at the TeV scale, and the low-energy value would be obtained by a small logarithmic evolution.  This would modify the numerical prediction for \(\alpha^{-1}\) and must be checked against the CODATA value.  Preliminary estimates suggest that the shift is within the current uncertainty of the topological prediction, but a precise calculation is required and will be presented in a future revision of Appendix~G.
	
	% ============================================================
	\section{Discussion: Is YUCT a Theory of Everything or of the Electroweak Scale?}
	\label{sec:discussion}
	
	The original YUCT framework aspired to explain all physical scales from the Planck mass down to the electron mass through a single coordination hierarchy.  The corrected scalar mass calculation reveals a tension within this ambition: if the coordination network operates at the Planck scale, the radiative corrections are enormous and must be cancelled by fine-tuned bare masses.
	
	Scenario B resolves this tension by reinterpreting the YUCT mass ladder.  Instead of producing masses directly through the depth \(L\), the ladder determines the \emph{Yukawa couplings} that, in turn, generate masses radiatively at the scale \(\LUV \sim 9\)~TeV.  The Planck scale then becomes the scale at which the coordination network \emph{begins} to form, but the physical granularity -- the scale at which the form factor becomes important -- is much lower.
	
	This reinterpretation preserves the algebraic beauty of YUCT while eliminating the fine-tuning problem.  It also transforms YUCT from a speculative theory of quantum gravity into a falsifiable model of electroweak physics, testable at current and future colliders.
	
	% ============================================================
	\section{Conclusion}
	\label{sec:conclusion}
	
	We have corrected a computational error in the scalar mass correction within YUCT, obtaining the exact analytic value \(C_0 = 3\sqrt{\pi}/(8\sqrt{2}) \approx 0.46999\).  With the Planck-scale identification \(\LUV = 3\pi\Mpl\), the radiative correction to the Higgs mass is \(\delta m_H \approx 1.57\times 10^{18}\)~GeV, requiring a fine-tuned cancellation against the bare mass.
	
	As an alternative, we have proposed Scenario B, in which the bare Lagrangian at the Planck scale is scale-invariant and all masses are generated radiatively.  Self-consistency then demands that the fundamental coordination scale be \(\LUV \approx 8.96\)~TeV, a value accessible to the LHC.  This scenario makes concrete, falsifiable predictions for high-energy scattering processes.
	
	The choice between Scenario A (aesthetically unsatisfying but consistent with the Planck-scale origin of YUCT) and Scenario B (predictive and testable, but requiring a reinterpretation of the YUCT mass ladder) cannot be made on theoretical grounds alone.  It must be decided by experiment.  We urge the LHC collaborations to examine high-mass Drell-Yan and dijet data for evidence of the form factor suppression predicted by Eq.~(\ref{eq:F}) with \(\LUV \sim 9\)~TeV.  A null result would rule out Scenario B and, while Scenario A remains logically possible, it would weaken the case for YUCT as a solution to the hierarchy problem.
	
	\section*{Data Availability}
	All calculations are self-contained; the Python code used for numerical evaluation is available in the YPSDC repository \url{https://github.com/Alexey-Yakushev-YUCT/YPSDC}.
	
	\begin{thebibliography}{99}
		\bibitem{AppendixL} A.\,V.\,Yakushev, ``Appendix L: Fractal Coordination Error Scaling,'' in \emph{YUCT}, Zenodo (2026).
		\bibitem{AppendixV} A.\,V.\,Yakushev, ``Appendix V: Time as Entropy of Coordination Processes,'' in \emph{YUCT}, Zenodo (2026).
		\bibitem{AppendixG} A.\,V.\,Yakushev, ``Appendix G: Gravity as Geometric Tension of the Coordination Network and Cyclic Cosmology in YUCT,'' in \emph{YUCT}, Zenodo (2026).
		\bibitem{AppendixM} A.\,V.\,Yakushev, ``Appendix M: YUCT Cosmology -- Inflation as a Coordination Phase Transition,'' in \emph{YUCT}, Zenodo (2026).
		\bibitem{AppendixLambda} A.\,V.\,Yakushev, ``Appendix \(\Lambda\): Resolution of the Hierarchy and Cosmological Constant Problems within YUCT,'' in \emph{YUCT}, Zenodo (2026).
		\bibitem{AppendixJ} A.\,V.\,Yakushev, ``Appendix J: Complete Derivation of Standard Model Flavor Parameters from YUCT Topological Offsets,'' in \emph{YUCT}, Zenodo (2026).
		\bibitem{AppendixFerra} A.\,V.\,Yakushev, ``Appendix Ferra1.2: Universal Coordination Constants for Ordered and Disordered Phases,'' in \emph{YUCT}, Zenodo (2026).
	\end{thebibliography}
	
\end{document}